% ==========================================================================
% CS3807 - Deep Learning Laboratory - Experiment 1
% Implementation of a Single Layer Perceptron for Binary Classification
%
% All numbers, tables and logs in this report were taken directly from the
% executed notebook "perceptron_lab.ipynb" (scikit-learn / numpy, seed = 42).
% Nothing below was made up, everything is copied from the actual run.
%
% FIGURES: All figures were exported by the notebook using
%   plt.savefig(path, dpi=600, format='pdf')
% with matplotlib rcParams set to font.family='serif',
% font.serif=['Liberation Serif','Times New Roman'], font.size=13, so every
% .pdf file is already 600 DPI and set in a Times-like serif font.
% Put the figures/ folder (already generated by the notebook) next to this
% .tex file and compile with pdflatex, no extra steps needed.
% ==========================================================================

\documentclass[11pt]{article}
\usepackage[margin=0.8in]{geometry}
\usepackage[T1]{fontenc}
\usepackage{mathptmx}                 % Times-family text & math font
\usepackage{amsmath,amssymb,booktabs,array,longtable}
\usepackage{graphicx}
\graphicspath{{figures/}}
\usepackage{caption}
\usepackage{hyperref}
\usepackage{enumitem}
\setlist{nosep}
\usepackage{xcolor}

\begin{document}

\begin{center}
{\Large\textbf{Shiv Nadar University Chennai}}
\end{center}

\renewcommand{\arraystretch}{1.25}

\begin{tabular}{|p{3.6cm}|p{7.6cm}|}
\hline
Degree \& Branch & B.Tech Artificial Intelligence \& Data Science, Semester V\\ \hline
Subject & CS3807 -- Deep Learning Laboratory\\ \hline
Name & R S kirithic\\ \hline
Section & AIDS B \\ \hline
Register Number & 24011101082\\ \hline
\end{tabular}

\vspace{4mm}
\begin{center}
{\large\textbf{Experiment 1}}\\
{\large\textbf{Implementation of a Single Layer Perceptron for Binary Classification}}
\end{center}

\section*{1. Objective}
The goal of this experiment is to build a Single
Layer Perceptron from scratch (no major libraries imported ) for binary classification. We have to use  the perceptron learning rule,  the activation function, plotting the learning process, and checking the classifier on a real dataset.

\section*{2. Tasks Performed}
\begin{enumerate}
\item Downloaded and explored the Banknote Authentication dataset.
\item Performed exploratory data analysis (EDA) on the dataset.
\item Preprocessed the dataset, normalized the features and split it into train and test sets.
\item Implemented a Single Layer Perceptron from scratch.
\item Trained the model using the perceptron learning rule and logged the training progress.
\item Evaluated the classifier using accuracy, precision, recall, F1-score and a confusion matrix.
\item Compared the Step and Sigmoid activation functions.
\item Studied the effect of different learning rates on convergence.
\item Compared the from-scratch implementation with scikit-learn's Perceptron.
\item Showed why a single layer perceptron cannot solve the XOR problem.
\item Studied the effect of feature normalization on convergence.
\end{enumerate}

\section*{3. Background Theory}
An artificial neuron takes one or more input features, multiplies them by weights, adds a bias, and
passes the result through an activation function to give the output.

Weighted sum:
\[
z = w^T x + b
\]

Prediction:
\[
\hat y = f(z)
\]

Weight update rule:
\[
w_{new} = w_{old} + \eta(y - \hat y)x
\]

Bias update:
\[
b_{new} = b_{old} + \eta(y - \hat y)
\]

 $\eta$ i: learning rate. 
 The Single Layer Perceptron is a supervised learning
 algorithm and it can only solve problems where the classes are linearly separable.

\section*{4. Activation Functions}
\subsection*{Step Activation Function}
The perceptron used in this experiment uses only the Step activation function:
\[
f(z) =
\begin{cases}
1, & z \geq 0 \\
0, & z < 0
\end{cases}
\]

This gives a hard 0 or 1 output and only works for linearly separable problems.

\begin{center}
\begin{tabular}{|l|l|l|}
\hline
\textbf{Activation} & \textbf{Output Range} & \textbf{Typical Use} \\ \hline
Step & \{0, 1\} & Classical Perceptron \\ \hline
Sigmoid & (0, 1) & Binary Classification Output Layer \\ \hline
Tanh & (-1, 1) & Hidden Layers (Traditional Networks) \\ \hline
ReLU & [0, $\infty$) & Hidden Layers in CNNs and MLPs \\ \hline
Leaky ReLU & (-$\infty$, $\infty$) & Avoiding Dying ReLU \\ \hline
Softmax & (0, 1) & Multi-class Classification Output \\ \hline
\end{tabular}
\end{center}

\section*{5. Dataset Description}
Dataset: Banknote Authentication Dataset (UCI Machine Learning Repository).

\begin{center}
\begin{tabular}{|l|l|}
\hline
Instances & 1372 \\ \hline
Features & 4 numerical features \\ \hline
Classes & 2 \\ \hline
Missing Values & None \\ \hline
Task & Binary Classification \\ \hline
\end{tabular}
\end{center}

Features: variance, skewness, curtosis, entropy of the wavelet transformed image.

Target class: 0 for authentic banknote, 1 for forged banknote.

\section*{6. Experimental Procedure and Results}

\subsection*{Task 1: Dataset Exploration}
I loaded the dataset with pandas and printed the first five rows, the dataset shape, the missing
value count and the descriptive statistics.

{\scriptsize
\begin{verbatim}  
First f ive samples:
   variance  skewness  curtosis  entropy  class
0   3.62160    8.6661   -2.8073 -0.44699      0
1   4.54590    8.1674   -2.4586 -1.46210      0
2   3.86600   -2.6383    1.9242  0.10645      0
3   3.45660    9.5228   -4.0112 -3.59440      0
4   0.32924   -4.4552    4.5718 -0.98880      0

Dataset dimensions: 1372 rows x 5 columns

Missing values per column:
variance    0
skewness    0
curtosis    0
entropy     0
class       0
\end{verbatim}
}

\textbf{Descriptive statistics:}

\begin{center}
{\small
\begin{tabular}{|l|c|c|c|c|}
\hline
\textbf{Stat} & \textbf{variance} & \textbf{skewness} & \textbf{curtosis} & \textbf{entropy} \\ \hline
mean & 0.4337 & 1.9224 & 1.3976 & -1.1917 \\ \hline
std  & 2.8428 & 5.8690 & 4.3100 & 2.1010 \\ \hline
min  & -7.0421 & -13.7731 & -5.2861 & -8.5482 \\ \hline
25\%  & -1.7730 & -1.7082 & -1.5750 & -2.4135 \\ \hline
50\%  & 0.4962 & 2.3197 & 0.6166 & -0.5867 \\ \hline
75\%  & 2.8215 & 6.8146 & 3.1793 & 0.3948 \\ \hline
max  & 6.8248 & 12.9516 & 17.9274 & 2.4495 \\ \hline
\end{tabular}
}
\end{center}

Class distribution: 762 good   notes (class 0) and 610 forged notes (class 1). The dataset is balanced, so we don't need an algorithm for class weighting  .

\subsection*{Task 2: Exploratory Data Analysis}

\begin{figure}[h!]
\centering
\includegraphics[width=0.75\textwidth]{01_feature_histograms.pdf}
\caption{Feature Histograms.}
\end{figure}


\begin{figure}[h!]
\centering
\includegraphics[width=0.55\textwidth]{02_correlation_heatmap.pdf}
\caption{Correlation Heatmap.}
\end{figure}



\begin{figure}[h!]
\centering
\includegraphics[width=0.85\textwidth]{16_pairplot_all_combinations.pdf}
\caption{Pairwise Scatter Plots of All Features, split by class.}
\end{figure}

\noindent\textbf{Observation:} the two classes form largely separate clusters along the
variance and skewness axes, so the problem looks close to linearly separable.

\begin{figure}[h!]
\centering
\includegraphics[width=0.75\textwidth]{04_boxplots.pdf}
\caption{Feature Boxplots by Class.}
\end{figure}

\noindent\textbf{Observation:} variance shows the clearest separation between the two classes
out of all four features, which matches what is seen in the scatter plots.

\subsection*{Task 3: Data Preprocessing}
All four features were normalized using \texttt{StandardScaler} (zero mean, unit variance), and the
data was split into 80\% training and 20\% testing, using stratified sampling to keep the class
balance in both sets.

{\scriptsize
\begin{verbatim}
Training set size: 1097 samples
Testing set size : 275 samples
Feature mean after scaling (train): [ 0.01  -0.008  0.003 -0.01 ]
Feature std  after scaling (train): [0.999 1.006 1.011 1.005]
\end{verbatim}
}

\subsection*{Task 4: Perceptron Implementation (From Scratch)}
The perceptron was implemented from scratch with the following parts:
\begin{itemize}
\item Weight initialization: zeros by default.
\item Bias initialization: zero by default.
\item Step activation function.
\item Forward propagation: $z = w^T x + b$, then apply the step function.
\item Learning rule: $w_{new} = w_{old} + \eta(y-\hat y)x$, $b_{new} = b_{old} + \eta(y-\hat y)$.
\end{itemize}

\subsection*{Task 5: Model Training}
The model was trained with learning rate 0.01 for 50 epochs. Each epoch printed the number of
misclassified samples, the updated weights and the updated bias.

\textbf{Epoch-wise training log (selected epochs, full log has 50 rows):}

\begin{center}
{\small
\begin{tabular}{|c|c|c|c|c|c|c|}
\hline
\textbf{Epoch} & \textbf{Errors} & \textbf{$w_1$} & \textbf{$w_2$} & \textbf{$w_3$} & \textbf{$w_4$} & \textbf{Bias} \\ \hline
1  & 59 & -0.0680 & -0.0938 & -0.0678 & -0.0110 & -0.03 \\ \hline
5  & 20 & -0.0991 & -0.1295 & -0.1037 &  0.0102 & -0.05 \\ \hline
10 & 20 & -0.1307 & -0.1541 & -0.1250 & -0.0060 & -0.06 \\ \hline
20 & 19 & -0.1636 & -0.1946 & -0.1613 & -0.0064 & -0.07 \\ \hline
30 & 17 & -0.1751 & -0.2043 & -0.1901 & -0.0196 & -0.08 \\ \hline
40 & 18 & -0.1948 & -0.2244 & -0.1853 & -0.0197 & -0.09 \\ \hline
50 & 18 & -0.1941 & -0.2358 & -0.2109 & -0.0195 & -0.10 \\ \hline
\end{tabular}
}
\end{center}

($w_1,w_2,w_3,w_4$ are the weights for variance, skewness, curtosis and entropy in that order.)

\begin{figure}[h!]
\centering
\includegraphics[width=0.65\textwidth]{05_training_error_vs_epoch.pdf}
\caption{Training Error vs. Epoch.}
\end{figure}

\noindent\textbf{Observation:} The  misclassifications initially drop a lot , from 59 down to around 15 to 20, and then kind of  oscillate for the rest of training instead of
reaching zero. The minimum reached was 12 errors, at epochs 23, 35 and 39, but the model never converges fully in 50 epochs, meaning the two classes are close to but not perfectly linearly
separable.

\begin{figure}[h!]
\centering
\includegraphics[width=0.65\textwidth]{06_weight_evolution.pdf}
\caption{Weight Evolution over Epochs.}
\end{figure}


\begin{figure}[h!]
\centering
\includegraphics[width=0.65\textwidth]{07_bias_evolution.pdf}
\caption{Bias Evolution over Epochs.}
\end{figure}

\noindent\textbf{Observation:} the bias moves from 0 down to about -0.10 in small steps over the
50 epochs and does not swing back and forth much, so it stabilizes along with the weights.

\subsection*{Task 6: Model Evaluation}
On the 275-sample test set:

{\scriptsize
\begin{verbatim}
Accuracy  : 0.9855
Precision : 0.9683
Recall    : 1.0000
F1-score  : 0.9839
Confusion Matrix:
[[149   4]
 [  0 122]]
\end{verbatim}
}

Out of 275 test samples, 149 authentic notes were correctly classified, 4 authentic notes were
wrongly classified as forged, no forged notes were missed, and 122 forged notes were correctly
classified.

\begin{figure}[h!]
\centering
\includegraphics[width=0.5\textwidth]{08_confusion_matrix.pdf}
\caption{Confusion Matrix (Test Set).}
\end{figure}

\noindent\textbf{Observation:} the model gets a perfect recall of 1.0, meaning it never misses a
forged note in the test set. All 4 mistakes are false positives, authentic notes marked as forged,
which is the safer kind of error to make in this problem.
Since it means that  we never let a forged one be marked as real one which is what the priority should be.
\subsection*{Task 7: Additional Task 1, Step vs Sigmoid Activation}

\begin{figure}[h!]
\centering
\includegraphics[width=0.85\textwidth]{17_step_vs_sigmoid.pdf}
\caption{Step vs. Sigmoid Activation Functions.}
\end{figure}

\noindent\textbf{Discussion:} The Step function gives a hard 0 or 1 output with a sudden jump at
$z=0$. Its derivative is zero everywhere except at that one point, so it does not carry any useful
gradient information. The Sigmoid function gives a smooth (0,1) output and is differentiable
everywhere, with derivative $f'(z) = f(z)(1-f(z))$.
.

\subsection*{Task 8: Learning Rate Comparison}
With zero weight initialization, the sign of $w \cdot x + b$, and therefore the whole sequence of
misclassifications, does not change with the learning rate. To make the comparison meaningful, all
three runs below start from the same small random initial weights and bias instead of zero.

\begin{center}
\begin{tabular}{|c|c|c|}
\hline
\textbf{Learning Rate} & \textbf{Final Misclassified (epoch 50)} & \textbf{Test Accuracy} \\ \hline
0.001 & 9  & 0.9891 \\ \hline
0.01  & 13 & 0.9855 \\ \hline
0.1   & 17 & 0.9818 \\ \hline
\end{tabular}
\end{center}

\begin{figure}[h!]
\centering
\includegraphics[width=0.65\textwidth]{09_learning_rate_comparison.pdf}
\caption{Learning Rate Comparison: Convergence Behaviour.}
\end{figure}

\noindent\textbf{Observation:} the smallest learning rate, 0.001, gives the fewest final
misclassifications and the best test accuracy. Larger learning rates make bigger jumps at each
update, which pushes the decision boundary past the minima more often near the end of training and
gives a slightly noisier, less accurate result here.
It also has a tendency to oscillate more .

\subsection*{Task 9 (Optional): Decision Boundary Visualization}
A second perceptron was trained using only two features, variance and skewness, so the decision
boundary could be drawn in 2D. Test accuracy with just these two features was 0.8364, lower than
the 0.9855 achieved using all four features.

\begin{figure}[h!]
\centering
\includegraphics[width=0.6\textwidth]{10_decision_boundary.pdf}
\caption{Decision Boundary (variance vs. skewness).}
\end{figure}

\noindent\textbf{Observation:} the perceptron learns one straight line that separates most of the
authentic and forged notes using only these two features, but some points from each class fall on
the wrong side, since two of the four features are being ignored here.

\subsection*{Task 10 (Optional): Comparison with Scikit-learn's Perceptron}

\begin{center}
\begin{tabular}{|l|c|c|}
\hline
\textbf{Metric} & \textbf{From-Scratch Perceptron} & \textbf{Scikit-learn Perceptron} \\ \hline
Accuracy  & 0.9855 & 0.9745 \\ \hline
Precision & 0.9683 & 0.9457 \\ \hline
Recall    & 1.0000 & 1.0000 \\ \hline
F1-score  & 0.9839 & 0.9721 \\ \hline
\end{tabular}
\end{center}

\noindent\textbf{Observation:} the from-scratch implementation slightly outperforms scikit-learn's
Perceptron on this run, though both reach the same perfect recall. The small gap is most likely
because of differences in default settings, such as \texttt{max\_iter} and the exact stopping rule,
rather than any real difference in the learning algorithm itself.

\subsection*{Task 11: Additional Task, Why a Single Layer Perceptron Cannot Solve XOR}
A perceptron was trained on the XOR truth table: $(0,0) \to 0$, $(0,1) \to 1$, $(1,0) \to 1$,
$(1,1) \to 0$.

{\scriptsize
\begin{verbatim}
Final misclassifications on XOR after 50 epochs: 4 out of 4 samples
\end{verbatim}
}

\begin{figure}[h!]
\centering
\includegraphics[width=0.65\textwidth]{11_xor_error_curve.pdf}
\caption{Perceptron on XOR: Error Never Reaches Zero.}
\end{figure}

\noindent\textbf{Observation:} the perceptron never learns XOR, it stays at 4 out of 4
misclassified even after 50 epochs. This i s because it has to fit a straight line to seperate it but no straight line can (impossible ) so we the two positive points $(0,1)$ and $(1,0)$ sit on opposite corners from each other
and from the two negative points. Separating them needs at least two lines, which means at least
one hidden layer, so XOR needs a multi-layer network to solve.
SO a non linear boundary for us to separate them .

\subsection*{Task 12: Additional Task, Effect of Feature Normalization on Convergence}

{\scriptsize
\begin{verbatim}
Without normalization -> epoch-1 errors: 63, mean errors (epochs 2-10): 23.1,
    mean errors (last 20 epochs): 15.2
With normalization    -> epoch-1 errors: 59, mean errors (epochs 2-10): 20.8,
    mean errors (last 20 epochs): 16.6
\end{verbatim}
}

\begin{figure}[h!]
\centering
\includegraphics[width=0.65\textwidth]{12_normalization_effect.pdf}
\caption{Effect of Feature Normalization on Convergence.}
\end{figure}

\noindent\textbf{Observation:} normalization gives slightly fewer errors early in training (epoch
1 and the epochs 2 to 10 average).However there isn't much difference .

\section*{7. Performance Tables}

\textbf{Training Summary}

\begin{center}
\begin{tabular}{|p{5.5cm}|p{6.5cm}|}
\hline
Dataset Size & 1372 \\ \hline
Train/Test Split & 80/20 (stratified), 1097 train, 275 test \\ \hline
Learning Rate & 0.01 \\ \hline
Epochs & 50 \\ \hline
Final Weights & [-0.1941, -0.2358, -0.2109, -0.0195] \\ \hline
Final Bias & -0.10 \\ \hline
Accuracy & 0.9855 \\ \hline
Precision & 0.9683 \\ \hline
Recall & 1.0000 \\ \hline
F1-score & 0.9839 \\ \hline
\end{tabular}
\end{center}

\section*{8. Discussion}

\begin{enumerate}
\item \textbf{Was the dataset linearly separable?}\\
yes but not fully . The training error drops fast in the first few epochs and then settles
between 12 and 20 misclassified samples out of 1097 training samples, and never reaches zero in 50
epochs. This means most of the data is linearly separable but a small number of points from the
two classes overlap.This is normal with real world data.

\item \textbf{Effect of learning rate.}\\
The smallest learning rate tried, 0.001, gave the best test accuracy (0.9891) and the fewest final
errors. Larger learning rates converge in a similar number of epochs but end up slightly less
accurate here, since bigger steps overshoot the boundary more near the end of training.

\item \textbf{From-scratch vs scikit-learn.}\\
Both implementations reached perfect recall (1.0) and very close accuracy, 0.9855 against 0.9745.
This is a good sign that the scratch implementation matches the standard algorithm correctly.

\item \textbf{Step vs Sigmoid.}\\
The Step function cannot be used in deep, multi-layer networks because it blocks gradient flow
during backpropagation. Differentiable activations like Sigmoid are needed for that.

\item \textbf{XOR.}\\
A single layer perceptron can only draw one straight decision boundary, so it cannot solve XOR,
which needs at least two lines to separate the classes. This is the classic limitation that shows
why multi-layer networks are needed, which will be covered in the next experiment.

\item \textbf{Normalization.}\\
Normalizing the features helped a bit in the early epochs of training but did not change the final
accuracy by much in this case, since the four raw features were not on wildly different scales to
begin with.
\end{enumerate}

\section*{9. Conclusion}
\begin{itemize}
\item A Single Layer Perceptron was implemented from scratch using the Rosenblatt learning rule
and the Step activation function, and applied to the Banknote Authentication dataset.
\item With learning rate 0.01 and 50 epochs, the model reached \textbf{98.55\% test accuracy},
with perfect recall and only 4 false positives out of 275 test samples.
\item Learning rate 0.001 gave the best result out of the three rates tried, reaching 98.91\% test
accuracy.
\item The from-scratch model performed close to, and slightly better than, scikit-learn's
Perceptron on this dataset and split.
\item The perceptron completely failed on the XOR problem, staying at 4 out of 4 misclassified
samples throughout training, which shows the basic limitation of a single linear decision boundary.
\item Feature normalization gave a small improvement early in training but did not change the
final accuracy much for this dataset.
\item Overall the experiment showed how a single artificial neuron learns through weight and bias
updates .
\end{itemize}

\section*{10. References}
\begin{enumerate}
\item F. Rosenblatt, ``The Perceptron,'' Psychological Review, 1958.
\item I. Goodfellow, Y. Bengio and A. Courville, \emph{Deep Learning}, MIT Press, 2016.
\item C. M. Bishop, \emph{Pattern Recognition and Machine Learning}, Springer, 2006.
\item S. Haykin, \emph{Neural Networks and Learning Machines}, Pearson, 2009.
\item UCI Machine Learning Repository -- Banknote Authentication Dataset.
\item Scikit-learn Documentation: Perceptron.
\end{enumerate}

\end{document}
